% --- Archon physics formalization source begin ---
% archon:physics
% archon:covers PhyXMiniProblems/problem_phyx_mini_0440.lean
% archon:source-report reports/phyx_mini/problem_phyx_mini_0440.source.json

\chapter{Physics problem phyx\_mini\_0440}
\label{ch:PhyXMiniProblems_problem_phyx_mini_0440}

\paragraph{Problem source.}
\#\# Physical scenario

A pressure cooker has the lid screwed on tight. A small opening with \$A = 5 \textbackslash{}, \textbackslash{}text\{mm\}\textasciicircum{}2\$ is covered with a petcock that can be lifted to let steam escape.

\#\# Question

How much mass should the petcock have to allow boiling at \$120 \textbackslash{}, \textasciicircum{}\{\textbackslash{}circ\}\textbackslash{}text\{C\}\$ with an outside atmosphere at \$101.3 \textbackslash{}, \textbackslash{}text\{kPa\}\$?

\#\# Answer choices

- A: A: 10 g

- B: B: 198.5 g

- C: C: 100 g

- D: D: 50 g

\#\# Figure caption (auxiliary; use the image as primary evidence)

The image depicts a schematic diagram of a closed container with a liquid and its vapor. 

- The container has a lid with a valve on top. Steam is shown exiting through the valve as indicated by wavy lines emerging from it.
- Inside the container, there are two main regions:
  - The lower region is labeled "Liquid," depicted with horizontal lines to represent the liquid phase.
  - The upper region is labeled "Steam or vapor," with dotted patterns to represent the gaseous phase filling the space above the liquid.
- The word "Steam" is labeled outside the container near the steam lines.

Overall, the image illustrates the phase equilibrium between a liquid and its vapor within a closed system.

\#\# Dataset metadata

Laws of Thermodynamics; Physical Model Grounding Reasoning, Multi-Formula Reasoning

\paragraph{Recorded answer/context.}
D: D: 50 g

\paragraph{Figure/image path.}
/root/proposal\_for\_physic/science-mango/phyx\_mini\_run/phyx\_data/test\_image/440.png

\paragraph{Formalization target.}
create a compiling Lean file with sorry bodies at `PhyXMiniProblems/problem\_phyx\_mini\_0440.lean`. The Lean declarations must preserve the physical quantities, dimensions or dimensional roles, figure labels, governing-law hypotheses, and final relation expressed by this problem.
Use Mathlib/Physlib names found through LeanExplore where available. If a physics API is missing, introduce faithful local abstractions rather than scalar placeholder aliases.

\begin{theorem}[Physics formalization target]
\label{thm:physics:phyx_mini_0440:target}
\lean{PhyXMiniProblems.ProblemPhyXMini0440.problem_phyx_mini_0440}
\uses{def:physics:phyx-mini-0440:phyxminiproblems-problemphyxmini0440-massinkilograms, def:physics:phyx-mini-0440:phyxminiproblems-problemphyxmini0440-massingrams, def:physics:phyx-mini-0440:phyxminiproblems-problemphyxmini0440-pressurecookersetup, def:physics:phyx-mini-0440:phyxminiproblems-problemphyxmini0440-matchespressurecookerscenario, def:physics:phyx-mini-0440:phyxminiproblems-problemphyxmini0440-matchessuppliedpressurecookerfigure, def:physics:phyx-mini-0440:phyxminiproblems-problemphyxmini0440-matchesproblemreadouts, def:physics:phyx-mini-0440:phyxminiproblems-problemphyxmini0440-usesstandardnearearthgravity, def:physics:phyx-mini-0440:phyxminiproblems-problemphyxmini0440-useswatersaturationtableat120c, def:physics:phyx-mini-0440:phyxminiproblems-problemphyxmini0440-satisfiesboilingandpetcockbalancelaws, def:physics:phyx-mini-0440:phyxminiproblems-problemphyxmini0440-answerchoice, def:physics:phyx-mini-0440:phyxminiproblems-problemphyxmini0440-isnearestdisplayedmass}
The assigned autoformalize agent should translate this physics problem into Lean declarations in the covered file, with theorem and lemma proof bodies written as `by sorry`.
\end{theorem}
\begin{proof}
This is an autoformalization task, not a proof task. Produce statements that can later be proved by the physics prover without weakening the physical model.
\end{proof}

% --- Archon declaration topology begin ---
\section{Declaration topology}
\begin{definition}[Mass Quantity]
  \label{def:physics:phyx-mini-0440:phyxminiproblems-problemphyxmini0440-massquantity}
  \lean{PhyXMiniProblems.ProblemPhyXMini0440.MassQuantity}
  A nonnegative physical mass, independent of the unit used to read it
\end{definition}
\begin{proof}
  This is the declared model definition of Mass Quantity.
\end{proof}
\begin{definition}[Acceleration Quantity]
  \label{def:physics:phyx-mini-0440:phyxminiproblems-problemphyxmini0440-accelerationquantity}
  \lean{PhyXMiniProblems.ProblemPhyXMini0440.AccelerationQuantity}
  A nonnegative acceleration magnitude with dimension L T⁻²
\end{definition}
\begin{proof}
  This is the declared model definition of Acceleration Quantity.
\end{proof}
\begin{definition}[Measured Temperature]
  \label{def:physics:phyx-mini-0440:phyxminiproblems-problemphyxmini0440-measuredtemperature}
  \lean{PhyXMiniProblems.ProblemPhyXMini0440.MeasuredTemperature}
  An absolute temperature and the zero-preserving unit in which its magnitude is stored. The affine Celsius offset is applied only by the scalar readout below, rather than treating Celsius as an absolute-temperature unit
\end{definition}
\begin{proof}
  This is the declared model definition of Measured Temperature.
\end{proof}
\begin{definition}[area Readout]
  \label{def:physics:phyx-mini-0440:phyxminiproblems-problemphyxmini0440-areareadout}
  \lean{PhyXMiniProblems.ProblemPhyXMini0440.areaReadout}
  Read an area in the square of a selected length unit
\end{definition}
\begin{proof}
  This is the declared model definition of area Readout.
\end{proof}
\begin{definition}[area In Square Meters]
  \label{def:physics:phyx-mini-0440:phyxminiproblems-problemphyxmini0440-areainsquaremeters}
  \lean{PhyXMiniProblems.ProblemPhyXMini0440.areaInSquareMeters}
  \uses{def:physics:phyx-mini-0440:phyxminiproblems-problemphyxmini0440-areareadout}
  Square-metre readout used in the pressure-force law
\end{definition}
\begin{proof}
  This is the declared model definition of area In Square Meters.
\end{proof}
\begin{definition}[area In Square Millimeters]
  \label{def:physics:phyx-mini-0440:phyxminiproblems-problemphyxmini0440-areainsquaremillimeters}
  \lean{PhyXMiniProblems.ProblemPhyXMini0440.areaInSquareMillimeters}
  \uses{def:physics:phyx-mini-0440:phyxminiproblems-problemphyxmini0440-areareadout}
  Square-millimetre readout used for the cooker opening
\end{definition}
\begin{proof}
  This is the declared model definition of area In Square Millimeters.
\end{proof}
\begin{definition}[mass Readout]
  \label{def:physics:phyx-mini-0440:phyxminiproblems-problemphyxmini0440-massreadout}
  \lean{PhyXMiniProblems.ProblemPhyXMini0440.massReadout}
  \uses{def:physics:phyx-mini-0440:phyxminiproblems-problemphyxmini0440-massquantity}
  Read a physical mass in a selected mass unit
\end{definition}
\begin{proof}
  This is the declared model definition of mass Readout.
\end{proof}
\begin{definition}[mass In Kilograms]
  \label{def:physics:phyx-mini-0440:phyxminiproblems-problemphyxmini0440-massinkilograms}
  \lean{PhyXMiniProblems.ProblemPhyXMini0440.massInKilograms}
  \uses{def:physics:phyx-mini-0440:phyxminiproblems-problemphyxmini0440-massquantity, def:physics:phyx-mini-0440:phyxminiproblems-problemphyxmini0440-massreadout}
  Kilogram readout used in the weight law
\end{definition}
\begin{proof}
  This is the declared model definition of mass In Kilograms.
\end{proof}
\begin{definition}[mass In Grams]
  \label{def:physics:phyx-mini-0440:phyxminiproblems-problemphyxmini0440-massingrams}
  \lean{PhyXMiniProblems.ProblemPhyXMini0440.massInGrams}
  \uses{def:physics:phyx-mini-0440:phyxminiproblems-problemphyxmini0440-massquantity, def:physics:phyx-mini-0440:phyxminiproblems-problemphyxmini0440-massreadout}
  Gram readout used by the displayed choices
\end{definition}
\begin{proof}
  This is the declared model definition of mass In Grams.
\end{proof}
\begin{definition}[acceleration In Meters Per Second Squared]
  \label{def:physics:phyx-mini-0440:phyxminiproblems-problemphyxmini0440-accelerationinmeterspersecondsquared}
  \lean{PhyXMiniProblems.ProblemPhyXMini0440.accelerationInMetersPerSecondSquared}
  \uses{def:physics:phyx-mini-0440:phyxminiproblems-problemphyxmini0440-accelerationquantity}
  Metre-per-second-squared readout of an acceleration
\end{definition}
\begin{proof}
  This is the declared model definition of acceleration In Meters Per Second Squared.
\end{proof}
\begin{definition}[pressure In Pascals]
  \label{def:physics:phyx-mini-0440:phyxminiproblems-problemphyxmini0440-pressureinpascals}
  \lean{PhyXMiniProblems.ProblemPhyXMini0440.pressureInPascals}
  Pascal readout of an absolute physical pressure
\end{definition}
\begin{proof}
  This is the declared model definition of pressure In Pascals.
\end{proof}
\begin{definition}[pressure In Kilopascals]
  \label{def:physics:phyx-mini-0440:phyxminiproblems-problemphyxmini0440-pressureinkilopascals}
  \lean{PhyXMiniProblems.ProblemPhyXMini0440.pressureInKilopascals}
  \uses{def:physics:phyx-mini-0440:phyxminiproblems-problemphyxmini0440-pressureinpascals}
  Kilopascal readout used by the atmospheric and steam-table data
\end{definition}
\begin{proof}
  This is the declared model definition of pressure In Kilopascals.
\end{proof}
\begin{definition}[temperature In Kelvin]
  \label{def:physics:phyx-mini-0440:phyxminiproblems-problemphyxmini0440-temperatureinkelvin}
  \lean{PhyXMiniProblems.ProblemPhyXMini0440.temperatureInKelvin}
  \uses{def:physics:phyx-mini-0440:phyxminiproblems-problemphyxmini0440-measuredtemperature}
  Kelvin readout of a measured absolute temperature
\end{definition}
\begin{proof}
  This is the declared model definition of temperature In Kelvin.
\end{proof}
\begin{definition}[temperature In Degrees Celsius]
  \label{def:physics:phyx-mini-0440:phyxminiproblems-problemphyxmini0440-temperatureindegreescelsius}
  \lean{PhyXMiniProblems.ProblemPhyXMini0440.temperatureInDegreesCelsius}
  \uses{def:physics:phyx-mini-0440:phyxminiproblems-problemphyxmini0440-measuredtemperature, def:physics:phyx-mini-0440:phyxminiproblems-problemphyxmini0440-temperatureinkelvin}
  Celsius readout, with 273.15 K represented exactly as 5463 / 20
\end{definition}
\begin{proof}
  This is the declared model definition of temperature In Degrees Celsius.
\end{proof}
\begin{definition}[pressure Force In Newtons]
  \label{def:physics:phyx-mini-0440:phyxminiproblems-problemphyxmini0440-pressureforceinnewtons}
  \lean{PhyXMiniProblems.ProblemPhyXMini0440.pressureForceInNewtons}
  \uses{def:physics:phyx-mini-0440:phyxminiproblems-problemphyxmini0440-areainsquaremeters, def:physics:phyx-mini-0440:phyxminiproblems-problemphyxmini0440-pressureinpascals}
  Normal force in newtons exerted by a uniform pressure on the opening
\end{definition}
\begin{proof}
  This is the declared model definition of pressure Force In Newtons.
\end{proof}
\begin{definition}[weight In Newtons]
  \label{def:physics:phyx-mini-0440:phyxminiproblems-problemphyxmini0440-weightinnewtons}
  \lean{PhyXMiniProblems.ProblemPhyXMini0440.weightInNewtons}
  \uses{def:physics:phyx-mini-0440:phyxminiproblems-problemphyxmini0440-massquantity, def:physics:phyx-mini-0440:phyxminiproblems-problemphyxmini0440-accelerationquantity, def:physics:phyx-mini-0440:phyxminiproblems-problemphyxmini0440-massinkilograms, def:physics:phyx-mini-0440:phyxminiproblems-problemphyxmini0440-accelerationinmeterspersecondsquared}
  Downward petcock weight in newtons
\end{definition}
\begin{proof}
  This is the declared model definition of weight In Newtons.
\end{proof}
\begin{definition}[Working Fluid]
  \label{def:physics:phyx-mini-0440:phyxminiproblems-problemphyxmini0440-workingfluid}
  \lean{PhyXMiniProblems.ProblemPhyXMini0440.WorkingFluid}
  The working liquid whose saturated vapor fills the cooker headspace
\end{definition}
\begin{proof}
  This is the declared model definition of Working Fluid.
\end{proof}
\begin{definition}[Petcock State]
  \label{def:physics:phyx-mini-0440:phyxminiproblems-problemphyxmini0440-petcockstate}
  \lean{PhyXMiniProblems.ProblemPhyXMini0440.PetcockState}
  Mechanical state of the weighted vent at the specified boiling point
\end{definition}
\begin{proof}
  This is the declared model definition of Petcock State.
\end{proof}
\begin{definition}[Figure Label]
  \label{def:physics:phyx-mini-0440:phyxminiproblems-problemphyxmini0440-figurelabel}
  \lean{PhyXMiniProblems.ProblemPhyXMini0440.FigureLabel}
  Literal labels visible in the supplied raster
\end{definition}
\begin{proof}
  This is the declared model definition of Figure Label.
\end{proof}
\begin{definition}[Figure Region]
  \label{def:physics:phyx-mini-0440:phyxminiproblems-problemphyxmini0440-figureregion}
  \lean{PhyXMiniProblems.ProblemPhyXMini0440.FigureRegion}
  Vertically ordered regions distinguished by the supplied raster
\end{definition}
\begin{proof}
  This is the declared model definition of Figure Region.
\end{proof}
\begin{definition}[Pressure Cooker Figure]
  \label{def:physics:phyx-mini-0440:phyxminiproblems-problemphyxmini0440-pressurecookerfigure}
  \lean{PhyXMiniProblems.ProblemPhyXMini0440.PressureCookerFigure}
  \uses{def:physics:phyx-mini-0440:phyxminiproblems-problemphyxmini0440-figurelabel, def:physics:phyx-mini-0440:phyxminiproblems-problemphyxmini0440-figureregion}
  Qualitative and geometric information read from the primary image. This record contains no numerical mass, pressure, area, or temperature datum
\end{definition}
\begin{proof}
  This is the declared model definition of Pressure Cooker Figure.
\end{proof}
\begin{definition}[Pressure Cooker Setup]
  \label{def:physics:phyx-mini-0440:phyxminiproblems-problemphyxmini0440-pressurecookersetup}
  \lean{PhyXMiniProblems.ProblemPhyXMini0440.PressureCookerSetup}
  \uses{def:physics:phyx-mini-0440:phyxminiproblems-problemphyxmini0440-massquantity, def:physics:phyx-mini-0440:phyxminiproblems-problemphyxmini0440-accelerationquantity, def:physics:phyx-mini-0440:phyxminiproblems-problemphyxmini0440-measuredtemperature, def:physics:phyx-mini-0440:phyxminiproblems-problemphyxmini0440-workingfluid, def:physics:phyx-mini-0440:phyxminiproblems-problemphyxmini0440-petcockstate, def:physics:phyx-mini-0440:phyxminiproblems-problemphyxmini0440-pressurecookerfigure}
  Independent physical quantities at the requested operating point. In particular, petcockMass is not defined using a displayed answer, and the saturation-pressure function represents a steam table independently of the actual cooker pressure
\end{definition}
\begin{proof}
  This is the declared model definition of Pressure Cooker Setup.
\end{proof}
\begin{definition}[Matches Pressure Cooker Scenario]
  \label{def:physics:phyx-mini-0440:phyxminiproblems-problemphyxmini0440-matchespressurecookerscenario}
  \lean{PhyXMiniProblems.ProblemPhyXMini0440.MatchesPressureCookerScenario}
  \uses{def:physics:phyx-mini-0440:phyxminiproblems-problemphyxmini0440-pressurecookersetup}
  Qualitative apparatus assumptions stated in the problem
\end{definition}
\begin{proof}
  This is the declared model definition of Matches Pressure Cooker Scenario.
\end{proof}
\begin{definition}[Matches Supplied Pressure Cooker Figure]
  \label{def:physics:phyx-mini-0440:phyxminiproblems-problemphyxmini0440-matchessuppliedpressurecookerfigure}
  \lean{PhyXMiniProblems.ProblemPhyXMini0440.MatchesSuppliedPressureCookerFigure}
  \uses{def:physics:phyx-mini-0440:phyxminiproblems-problemphyxmini0440-pressurecookersetup}
  Primary-image evidence: liquid occupies the lower region, steam or vapor the upper region, and escaping steam is drawn above the lid-mounted petcock
\end{definition}
\begin{proof}
  This is the declared model definition of Matches Supplied Pressure Cooker Figure.
\end{proof}
\begin{definition}[Matches Problem Readouts]
  \label{def:physics:phyx-mini-0440:phyxminiproblems-problemphyxmini0440-matchesproblemreadouts}
  \lean{PhyXMiniProblems.ProblemPhyXMini0440.MatchesProblemReadouts}
  \uses{def:physics:phyx-mini-0440:phyxminiproblems-problemphyxmini0440-areainsquaremillimeters, def:physics:phyx-mini-0440:phyxminiproblems-problemphyxmini0440-pressureinkilopascals, def:physics:phyx-mini-0440:phyxminiproblems-problemphyxmini0440-temperatureindegreescelsius, def:physics:phyx-mini-0440:phyxminiproblems-problemphyxmini0440-pressurecookersetup}
  Numerical data stated in the prose. No petcock-mass value or answer label occurs here
\end{definition}
\begin{proof}
  This is the declared model definition of Matches Problem Readouts.
\end{proof}
\begin{definition}[Uses Standard Near Earth Gravity]
  \label{def:physics:phyx-mini-0440:phyxminiproblems-problemphyxmini0440-usesstandardnearearthgravity}
  \lean{PhyXMiniProblems.ProblemPhyXMini0440.UsesStandardNearEarthGravity}
  \uses{def:physics:phyx-mini-0440:phyxminiproblems-problemphyxmini0440-accelerationinmeterspersecondsquared, def:physics:phyx-mini-0440:phyxminiproblems-problemphyxmini0440-pressurecookersetup}
  Standard near-Earth gravity used by the recorded numerical choice. It is an independent environmental calibration rather than the requested mass
\end{definition}
\begin{proof}
  This is the declared model definition of Uses Standard Near Earth Gravity.
\end{proof}
\begin{definition}[Uses Water Saturation Table At120 C]
  \label{def:physics:phyx-mini-0440:phyxminiproblems-problemphyxmini0440-useswatersaturationtableat120c}
  \lean{PhyXMiniProblems.ProblemPhyXMini0440.UsesWaterSaturationTableAt120C}
  \uses{def:physics:phyx-mini-0440:phyxminiproblems-problemphyxmini0440-pressureinkilopascals, def:physics:phyx-mini-0440:phyxminiproblems-problemphyxmini0440-pressurecookersetup}
  The saturated-water steam-table readout at the independently specified 120 °C operating temperature is 198.5 kPa. This thermodynamic datum does not constrain the petcock mass
\end{definition}
\begin{proof}
  This is the declared model definition of Uses Water Saturation Table At120 C.
\end{proof}
\begin{definition}[Satisfies Boiling And Petcock Balance Laws]
  \label{def:physics:phyx-mini-0440:phyxminiproblems-problemphyxmini0440-satisfiesboilingandpetcockbalancelaws}
  \lean{PhyXMiniProblems.ProblemPhyXMini0440.SatisfiesBoilingAndPetcockBalanceLaws}
  \uses{def:physics:phyx-mini-0440:phyxminiproblems-problemphyxmini0440-pressureforceinnewtons, def:physics:phyx-mini-0440:phyxminiproblems-problemphyxmini0440-weightinnewtons, def:physics:phyx-mini-0440:phyxminiproblems-problemphyxmini0440-pressurecookersetup}
  The governing laws are separate from their numerical calibrations: water boils when the headspace steam pressure equals its saturation pressure at the water temperature; at incipient lift, upward steam-pressure force balances downward atmospheric-pressure force and petcock weight. Neither law asserts a numerical petcock mass or selects an answer choice
\end{definition}
\begin{proof}
  This is the declared model definition of Satisfies Boiling And Petcock Balance Laws.
\end{proof}
\begin{definition}[Answer Choice]
  \label{def:physics:phyx-mini-0440:phyxminiproblems-problemphyxmini0440-answerchoice}
  \lean{PhyXMiniProblems.ProblemPhyXMini0440.AnswerChoice}
  Labels of the four answers printed with the problem
\end{definition}
\begin{proof}
  This is the declared model definition of Answer Choice.
\end{proof}
\begin{definition}[displayed Mass In Grams]
  \label{def:physics:phyx-mini-0440:phyxminiproblems-problemphyxmini0440-displayedmassingrams}
  \lean{PhyXMiniProblems.ProblemPhyXMini0440.displayedMassInGrams}
  \uses{def:physics:phyx-mini-0440:phyxminiproblems-problemphyxmini0440-answerchoice}
  Gram values printed beside the four labels. This records the candidate list without asserting which candidate follows from the physical laws
\end{definition}
\begin{proof}
  This is the declared model definition of displayed Mass In Grams.
\end{proof}
\begin{definition}[Is Nearest Displayed Mass]
  \label{def:physics:phyx-mini-0440:phyxminiproblems-problemphyxmini0440-isnearestdisplayedmass}
  \lean{PhyXMiniProblems.ProblemPhyXMini0440.IsNearestDisplayedMass}
  \uses{def:physics:phyx-mini-0440:phyxminiproblems-problemphyxmini0440-massingrams, def:physics:phyx-mini-0440:phyxminiproblems-problemphyxmini0440-pressurecookersetup, def:physics:phyx-mini-0440:phyxminiproblems-problemphyxmini0440-answerchoice, def:physics:phyx-mini-0440:phyxminiproblems-problemphyxmini0440-displayedmassingrams}
  A displayed mass choice is at least as close as every alternative
\end{definition}
\begin{proof}
  This is the declared model definition of Is Nearest Displayed Mass.
\end{proof}
% --- Archon declaration topology end ---
% --- Archon physics formalization source end ---
